\documentclass[11pt,a4paper]{article}
\usepackage[utf8]{inputenc}
\usepackage[T1]{fontenc}
\usepackage[english]{babel}
\usepackage{amsmath, amsthm, amssymb, amsfonts}
\usepackage{enumerate}
\usepackage{graphicx}
\usepackage{fancyhdr}
\usepackage{booktabs}
\usepackage{float}
\usepackage[colorlinks=true,breaklinks=false]{hyperref}
\usepackage{listings}
\usepackage[norelsize,ruled,vlined]{algorithm2e}
\usepackage{csquotes}
\usepackage{url}
\usepackage{ifthen}
\usepackage{keyval}
\usepackage{etoolbox}
\usepackage{etex}
\usepackage{datetime}
\usepackage{eurosym}
\usepackage{tikz}
\usepackage{wrapfig}
\usepackage{paralist}
\usepackage[backend=bibtex]{biblatex}
\usepackage{xcolor}
\usepackage{caption}
\DeclareCaptionFont{white}{\color{white}}
\DeclareCaptionFormat{listing}{\colorbox[cmyk]{0.43, 0.35, 0.35,0.01}{\parbox{\textwidth}{\hspace{15pt}#1#2#3} } }
\captionsetup[lstlisting]{format=listing, labelfont=white, textfont=white, singlelinecheck=false, margin=0pt, font={bf,footnotesize} }

\addbibresource{project.bib}

% Single spacing after dots and colons
\frenchspacing
% \hyphenation{directembedding}
% \hyphenation{DSLVirtualization}
% \hyphenation{ReificationTransformation}

%%%%%%%%%%%%%%%%%%%%%%%%%%%%%%%%%%%%%%%%%%%%%%%%%%%%%%%%%%%%%
%                        Setup
%%%%%%%%%%%%%%%%%%%%%%%%%%%%%%%%%%%%%%%%%%%%%%%%%%%%%%%%%%%%%

\begin{document}

\definecolor{dkgreen}{rgb}{0,0.6,0}
\definecolor{gray}{rgb}{0.5,0.5,0.5}
\definecolor{mauve}{rgb}{0.58,0,0.82}

\lstset{language=Scala,
  aboveskip=3mm,
  basicstyle={\small\ttfamily},
  belowskip=3mm,
  breakatwhitespace=true,
  breaklines=true,
  columns=flexible,
  commentstyle=\color{dkgreen},
  keywordstyle=\color{blue},
  numbers=left,
  numberstyle=\tiny\color{gray},
  showstringspaces=false,
  stepnumber=1,
  stringstyle=\color{mauve},
  tabsize=2,
}

\usetikzlibrary{calc,trees,positioning,arrows,chains,shapes.geometric,%
    decorations.pathreplacing,decorations.pathmorphing,shapes,%
    matrix,shapes.symbols}
\tikzset{>=stealth',
  punktchain/.style={rectangle,
    rounded corners,
    % fill=black!10,
    draw=black, very thick,
    text width=12.5em,
    minimum height=2.5em,
    text centered,
    on chain},
  every join/.style={->, shorten >=2pt},
}


% Example of title page for the projects carried out within the lasec

% Simply include it in your mastex tex file:
%        % Example of title page for the projects carried out within the lasec

% Simply include it in your mastex tex file:
%        % Example of title page for the projects carried out within the lasec

% Simply include it in your mastex tex file:
%        \input{cover}


% Updated March 2006 (SP)


\newcommand{\logoepfl}[0]{
  \begin{center}
    \includegraphics[width=4cm]{logo_epfl_coul.eps}
  \end{center}
  \vspace{0.3cm}
  \hrule
}
\newcommand{\logolasec}[0]{
  \vspace{1cm}
  \hrule
  \begin{center}
    \includegraphics[width=4.5cm]{logo_lasec_coul.eps}
  \end{center}
}
\newcommand{\project}[1]{
  \begin{center}
    \large{#1}
  \end{center}
  \vspace{1cm}
}
\newcommand{\department}[1]{
  \begin{center}
    \large{#1}
  \end{center}
}
\newcommand{\supervisor}[3]{
  \begin{center}
    \begin{normalsize}{
        \bf #1}\\#2\\#3
    \end{normalsize}
  \end{center}
}
\renewcommand{\author}[1]{
  \begin{center}
    \Large{#1}
  \end{center}
  \vspace{0.5cm}
}
\renewcommand{\title}[1]{
  \vspace{3cm}
  \begin{center}
    \huge{#1}
  \end{center}
  \vspace{1.7cm}
}
\renewcommand{\date}[2]{
  \begin{center}
    \normalsize{#1 #2}
  \end{center}
  \vspace{0.5cm}
}


\thispagestyle{empty}


% begin title page
  \logoepfl

  \title{scalafmt: yet another approach \\ to code formatting}

  \author{Ólafur Páll Geirsson}
  \department{School of Computer and Communication Sciences}
  \project{Master's Thesis}

  \date{June}{2015}

  \begin{center}
    \begin{tabular}{cc}
      \begin{tabular}{p{4.0cm}}
        \supervisor{Responsible}{Prof. Martin Odersky}{EPFL / LAMP}
      \end{tabular}&
      \begin{tabular}{p{4.0cm}}
        \supervisor{Supervisor}{Eugene Burmako}{EPFL / LAMP}
      \end{tabular}
    \end{tabular}
  \end{center}

% end title page




% Updated March 2006 (SP)


\newcommand{\logoepfl}[0]{
  \begin{center}
    \includegraphics[width=4cm]{logo_epfl_coul.eps}
  \end{center}
  \vspace{0.3cm}
  \hrule
}
\newcommand{\logolasec}[0]{
  \vspace{1cm}
  \hrule
  \begin{center}
    \includegraphics[width=4.5cm]{logo_lasec_coul.eps}
  \end{center}
}
\newcommand{\project}[1]{
  \begin{center}
    \large{#1}
  \end{center}
  \vspace{1cm}
}
\newcommand{\department}[1]{
  \begin{center}
    \large{#1}
  \end{center}
}
\newcommand{\supervisor}[3]{
  \begin{center}
    \begin{normalsize}{
        \bf #1}\\#2\\#3
    \end{normalsize}
  \end{center}
}
\renewcommand{\author}[1]{
  \begin{center}
    \Large{#1}
  \end{center}
  \vspace{0.5cm}
}
\renewcommand{\title}[1]{
  \vspace{3cm}
  \begin{center}
    \huge{#1}
  \end{center}
  \vspace{1.7cm}
}
\renewcommand{\date}[2]{
  \begin{center}
    \normalsize{#1 #2}
  \end{center}
  \vspace{0.5cm}
}


\thispagestyle{empty}


% begin title page
  \logoepfl

  \title{scalafmt: yet another approach \\ to code formatting}

  \author{Ólafur Páll Geirsson}
  \department{School of Computer and Communication Sciences}
  \project{Master's Thesis}

  \date{June}{2015}

  \begin{center}
    \begin{tabular}{cc}
      \begin{tabular}{p{4.0cm}}
        \supervisor{Responsible}{Prof. Martin Odersky}{EPFL / LAMP}
      \end{tabular}&
      \begin{tabular}{p{4.0cm}}
        \supervisor{Supervisor}{Eugene Burmako}{EPFL / LAMP}
      \end{tabular}
    \end{tabular}
  \end{center}

% end title page




% Updated March 2006 (SP)


\newcommand{\logoepfl}[0]{
  \begin{center}
    \includegraphics[width=4cm]{logo_epfl_coul.eps}
  \end{center}
  \vspace{0.3cm}
  \hrule
}
\newcommand{\logolasec}[0]{
  \vspace{1cm}
  \hrule
  \begin{center}
    \includegraphics[width=4.5cm]{logo_lasec_coul.eps}
  \end{center}
}
\newcommand{\project}[1]{
  \begin{center}
    \large{#1}
  \end{center}
  \vspace{1cm}
}
\newcommand{\department}[1]{
  \begin{center}
    \large{#1}
  \end{center}
}
\newcommand{\supervisor}[3]{
  \begin{center}
    \begin{normalsize}{
        \bf #1}\\#2\\#3
    \end{normalsize}
  \end{center}
}
\renewcommand{\author}[1]{
  \begin{center}
    \Large{#1}
  \end{center}
  \vspace{0.5cm}
}
\renewcommand{\title}[1]{
  \vspace{3cm}
  \begin{center}
    \huge{#1}
  \end{center}
  \vspace{1.7cm}
}
\renewcommand{\date}[2]{
  \begin{center}
    \normalsize{#1 #2}
  \end{center}
  \vspace{0.5cm}
}


\thispagestyle{empty}


% begin title page
  \logoepfl

  \title{scalafmt: yet another approach \\ to code formatting}

  \author{Ólafur Páll Geirsson}
  \department{School of Computer and Communication Sciences}
  \project{Master's Thesis}

  \date{June}{2015}

  \begin{center}
    \begin{tabular}{cc}
      \begin{tabular}{p{4.0cm}}
        \supervisor{Responsible}{Prof. Martin Odersky}{EPFL / LAMP}
      \end{tabular}&
      \begin{tabular}{p{4.0cm}}
        \supervisor{Supervisor}{Eugene Burmako}{EPFL / LAMP}
      \end{tabular}
    \end{tabular}
  \end{center}

% end title page



\begin{abstract}
  Code formatters bring many benefits to software development such as enforcing a consistent coding style across teams, more effective code reviews and enabling automated large-scale refactoring.
  % Still, code formatters can be tricky to get right.
  This thesis addresses how to develop a code formatter for the Scala programming language.
  We present scalafmt, an opinionated Scala code formatter that captures many popular Scala idioms and coding styles.
  This thesis introduces language-agnostic algorithms and tooling that scalafmt uses to implement advanced features such as line wrapping and configurable vertical alignment.
  We have validated that these techniques work well in practice.
  Scalafmt has been installed over 6.500 times in only 3 months and
  several popular open-source libraries have chosen to reformat their codebases with scalafmt.
\end{abstract}
% \newpage
\vspace{1in}
\renewcommand{\abstractname}{Útdráttur}

\begin{abstract}
  Kóðasniðlar (e. code formatters) eru nytsamleg tól í hugbúnaðarþróun.
  Helstu kostir kóðasniðla eru meðal annars að geta sjálfvirkt framfylgt samræmdum kóðastíl, skilvirkari kóðaumsögnum og gera kleift að endurskipleggja stór forritasöfn.
  Þetta verkefni fjallar um að þróa kóðasniðil fyrir Scala forritunarmálið.
  Við kynnum scalafmt, kóðasniðil fyrir Scala sem fangar marga vinsæla Scala kóðastíla og forritunartiltæki.
  Þetta verkefni kynnir reiknirit og gagnagrindur til að útfæra háþróaða eiginlega eins og að brjóta langar forritunarskipanir á einni línu niður í margar línur og raða tóka af svipuðu tagi frá mörgum línum þannig að tókarnir liggi á sama lóðrétta dálki.
  Aðferðir sem kynntar eru í þessu verkefni hafa sannað sig í verki.
  Scalafmt hefur verið halað niður yfir 6.500 sinnum á eingöngu 3 mánuðum og fjöldi af vinsælum opnum forritasöfnum hafa kosið að sníða kóðann sinn með scalafmt.
\end{abstract}

\begin{abstract}
  Automatic code formatters bring many benefits to software development, yet they can be tricky to get right.
  This thesis addresses the problem of developing a code formatter for the
  Scala programming language that captures the language's most popular idioms
  and coding styles.
  Our work has been limited to formatting Scala code.
  Still, we have developed data structures, algorithms and tools that we
  believe can applied to develop code formatters for a variety of other
  programming languages.
\end{abstract}
\tableofcontents
% \section{Introduction} % (fold)
\label{sec:Introduction}
\lstset{style=scala}
Without code formatters, software developers are responsible for manipulating all syntactic trivia in their programs.
What is syntactic trivia?
Consider the Scala code snippets in listings~\ref{lst:unformatted} and~\ref{lst:formatted1}.

\begin{minipage}{.45\textwidth}
\lstinputlisting[label={lst:unformatted}, caption=Unformatted code]{code/unformatted.scala}
\end{minipage}
\hfil
\begin{minipage}{.45\textwidth}
\lstinputlisting[label={lst:formatted1}, caption=Formatted code]{target/formatted1.scala}
\end{minipage}

Both snippets represent the same application.
The only difference lies in where the programmer has chosen to break lines.
Characters such as spaces and line breaks that do not affect the execution of the program are syntactic trivia.
Although syntactic trivia has no meaning for the execution of the program, listing~\ref{lst:formatted1} is arguably easier to understand, maintain and extend for the software developer.
A code formatter is a tool that automatically converts a program such as in listing~\ref{lst:unformatted} into a readable and maintainable program such as in listing~\ref{lst:formatted1}.
Code formatting brings several benefits to software development.

Code formatting enables automated large-scale refactoring.
Google used ClangFormat\autocite{jasper_clangformat_2013}, a C++ code formatter,
in the process of migrating a large legacy C++98 codebase to use the modern C++11 standard\autocite{wright_large-scale_2013}.
The automatically refactored code was formatted with ClangFormat to ensure that it adhered to Google's strict coding style\autocite{_google_????}.
Similar migrations can be expected in the near future for the Scala community once new dialects, such as Dotty\autocite{rompf_f_2015}, gain popularity.

Code formatting is valuable in collaborative coding environments.
The Scala.js project\autocite{_scala.js_????} has over 40 contributors and the Scala.js coding style\autocite{doeraene_scala.js_2015} --- which each Scala.js contributor is expected to know by heart --- is defined at a whopping 2.600 word count.
Each contributed patch is manually verified against the coding style by the project maintainers.
This adds a burden on both contributors and maintainers.
Several prominent Scala community members have raised this issue.
ENSIME\autocite{_ensime_????} is a popular Scala interaction mode for text editors such as Vim and Emacs.
Sam Halliday, an ENSIME maintainer, says ``I don't have time to talk about formatting in code reviews. I want the machine to do it so I can focus on the design.''\autocite{halliday_i_2016-1}.
Akka\autocite{_akka_????} is a popular concurrent and distributed programming library for Scala with over 300 contributors.
Viktor Klang, a maintainer of Akka, suggests a better alternative ``Code style should not be enforced by review, but by automate rewriting. Evolve the style using PRs against the rewriting config.''.\autocite{klang_code_2016}.

With scalafmt, we hope to relieve Scala developers from the burden of manipulating syntactic trivia so they can instead direct their full attention to writing correct, maintainable and fast code.

\subsection{Research objective}
% Code formatters bring several benefits to software development.
% However, the Scala community has enjoyed limited benefits from the 
% Can we develop a code formatter for the Scala programming language

What algorithms and data structures allow us to develop a code formatter for the Scala programming language with the following properties?
\begin{itemize}
  \item \textbf{Maximum line length setting}: a code formatter with a maximum line length setting ensures that each line in the formatted output contains no more than a certain number of characters.
    Many coding styles enforce a maximum line length in order to make code readable from different environments such as small split screens and code review interfaces.
  \item \textbf{Opinionated}:
    An opinionated setting is a prerequisite to enforce a uniform coding style.
    An opinionated code formatter takes liberty to disregard line breaks and other formatting decisions in the original source input to ensure that formatted source files follow the same line breaking conventions.
  \item \textbf{Vertical alignment}: vertical alignment is a formatting convention where redundant whitespace is added before a token to align it on the same vertical column as similar tokens from other lines. Many Scala coding styles enforce vertical alignment to enhance code readability.
  \item \textbf{Performance}: Users expect code formatters to run fast.
    In the most demanding settings, a code formatter needs to be able to format thousands of lines of code in at most a few hundred milliseconds.

\end{itemize}

\subsection{Contributions}
The main contributions presented in this thesis are the following:
\begin{itemize}
  \item language-agnostic algorithms and data structures to implement line
    wrapping with a maximum line length setting --- in our opinion, the most
    challenging part of developing a code formatter --- as well as configurable
    vertical alignment.
    This work is presented in section~\ref{sec:algorithms}.
  \item methods to optimize and test the developed algorithms.
    This work is presented in section~\ref{sec:tooling}.
  \item practical validation of developed algorithms and methods with the scalafmt code formatter.
    The empirical results of this validation are presented in section~\ref{sec:evaluation}.
\end{itemize}

\section{Introduction} % (fold)
\label{sec:Introduction}
Software developers use code formatters to make their code look pretty.
A code formatter should not change the behaviour of code.
Typically, a code formatter only modifies whitespace tokens, or syntactic trivia.
Syntactic trivia has nothing to do with the behavior of code.
Code formatters promise to relieve software developers from manually manipulating
syntactic trivia.
Instead, developers can focus on writing correct, performant and maintainable code.

Code formatters brings many other benefits to software development.
They can help enforce a consistent coding style in a codebase that's touched by multiple developers.
A consistent code style aids readability of code.
Code formatting also enables automated refactoring tools, see~\autocite{wright_large-scale_2013}.
Code formatting also lowers the barrier to entry for novice programmers, who are not aware of the coding style conventions.

Code formatters have existed for a long time.
In more recent years we've seen many new opimization based code formatters, such as clang-format and dartfmt.
Optimization based code formatters define a fit function for a given layout.
A good layout has low cost while a bad layout has a high cost.
The formatter runs an algorithm to choose the layout with the lowest code.
This idea is not new, for example TeX uses a fit function to optimally break lines
in text documents.

Much work has been put into developing language agnostic code formatters or pretty printers.
Researchers of such work correctly highlights that developing a custom code formatter for each language requires a lot of effort.
It may be possible that we find a perfect way to develop language independent code formatters, see~\autocite{mps_article}.
Still, today no such framework exists and people who write Scala code each day are still choosing to manually format their code.
The goal of this thesis is not to address language agnostic code formatting.

\subsection{Problem statement}
This thesis addresses the problem of developing a code formatter, \texttt{scalafmt}, for the Scala programming language.
This formatter should:
\begin{itemize}
  \item capture Scala's most common idioms and popular coding styles, most importanly a line length limit.
  \item be opinionated, it should be able to produce nice looking code even from computer generated input.
  \item be fast, <200ms for 99 percentile of source files, <2s for rest.
\end{itemize}

\subsection{Contributions}
The main contributions presented in this thesis are the following:
\begin{itemize}
    \item We have developed a framework of data structures, algorithms and tooling that help
      developing code formatters. This work is presented in section X.
    \item We have applied this framework to develop \texttt{scalafmt}, a code
      formatter for the Scala programming language.
\end{itemize}

\section{Background}
\subsection{Scala the programming language}

Throughout the paper we assume familiarity with the basics of the Scala
Programming Language~\autocite{odersky_scala_2004}.
\subsubsection{Higher order functions}
\subsubsection{Pattern matching}
\subsubsection{Metaprogramming with scala.meta}
\begin{itemize}
  \item Tree nodes have parent links.
  \item Token classes are types.
\end{itemize}
\subsection{Code formatting}
\subsubsection{Scalariform}
\section{Framework}
\subsection{Data structures}
\subsubsection{FormatToken}
\subsubsection{Split combinator}
\subsubsection{Policy combinator}
\subsubsection{OptimalToken combinator}
\subsubsection{State}
\subsubsection{Indent}
\subsection{Algorithms}
Pipeline of three stages: pre-processing, line wrapping and post-processing.
\subsubsection{Pre-processing}
 \begin{itemize}
   \item AST vs. token stream
 \end{itemize}
rustfmt considers that all eventually converge to a hybrid of the two.
Indeed, clang-format has an elaborate parser.
\subsubsection{Layout}
\begin{itemize}
  \item Router
  \item Search.
\end{itemize}
\subsubsection{Post-processing}
\begin{itemize}
  \item vertical alignment
  \item comment formatting
  \item stripMargin alignment
\end{itemize}
\subsection{Tooling}
\subsubsection{Heatmaps}
\subsubsection{Unit tests}
\subsubsection{Property based testing}
\begin{enumerate}
  \item AST integrity
  \item Idempotent
\end{enumerate}
\section{scalafmt}
\subsection{Custom optimizations}
\begin{enumerate}
  \item dequeueOnNewStatements
  \item recurseOnBlocks
  \item escapeInPathologicalCases
  \item pruneSlowStates
\end{enumerate}
\section{Evaluation}
\subsection{Performance}
\subsection{Output}
\begin{itemize}
  \item gofmt
  \item rustfmt
  \item dartfmt
  \item clang-format
  \item scalariform
\end{itemize}
\section{Related work}
\subsection{Combinator based}
\begin{enumerate}
  \item Houghes 1995
  \item Wadler 1999
\end{enumerate}
\subsection{Optimization-oriented}
\begin{enumerate}
  \item clang-format Dijkstra's 2010
  \item dartfmt Best-first search 2014
  \item rfmt 2015
\end{enumerate}
\begin{enumerate}
  \item Optimal line breaking
  \item Oppen
\end{enumerate}
\section{Discussion}
\subsection{Future work}
\subsection{Conclusion}
\printbibliography{}

\end{document}

